\PassOptionsToPackage{unicode}{hyperref}
\PassOptionsToPackage{hyphens}{url}
\documentclass[9pt,twocolumn]{extarticle}
\usepackage[letterpaper, top=1in, bottom=1in, left=0.75in, right=0.75in]{geometry}
\usepackage{times}
\usepackage{mathptmx}
\usepackage{helvet}
\renewcommand{\familydefault}{\rmdefault}
\usepackage{xcolor}
\usepackage{amsmath,amssymb}
\usepackage{graphicx}
\usepackage{booktabs}
\usepackage{array}
\usepackage{multirow}
\usepackage{calc}
\usepackage{etoolbox}
\usepackage{url}
\usepackage{hyperref}
\usepackage{bookmark}
\usepackage{setspace}
\usepackage{caption}
\captionsetup{font=small,labelfont=bf}
\usepackage{enumitem}
\setlist{nosep,leftmargin=*}
\setlength{\parindent}{0pt}
\setlength{\parskip}{3pt plus 1pt minus 1pt}
\setlength{\emergencystretch}{3em}
\providecommand{\tightlist}{\setlength{\itemsep}{0pt}\setlength{\parskip}{0pt}}
\urlstyle{same}
\hypersetup{hidelinks,pdfcreator={LaTeX via pandoc}}

\title{\textbf{\sffamily\fontsize{18}{22}\selectfont VeritasCarbon: Traceable Instruction Data Generation for ESG}}
\author{%
\normalsize
\textbf{Yihan Jiang}\\ Society Hub, HKUST(GZ)\\ yjiang909@connect.hkust-gz.edu.cn
\and
\textbf{Fei Peng}\\ Nanchang University\\ pengfei24@email.ncu.edu.cn
\and
\textbf{Kok Sin Woon}\\ Society Hub, HKUST(GZ)\\ koksinwoon@hkust-gz.edu.cn
\and
\textbf{Qianping Ren}\\ Society Hub, HKUST(GZ)\\ qianpingren@hkust-gz.edu.cn
\and
\textbf{Yichang Xu}\\ Sun Yat-sen University\\ xuych73@mail2.sysu.edu.cn
\and
\textbf{Yujing Yang}\\ Sun Yat-sen University\\ Yangyj256@mail2.sysu.edu.cn
}
\date{\vspace{-2ex}}

\begin{document}
\maketitle
\begin{abstract}
Fine-tuning large language models for ESG (Environmental, Social, and
Governance) tasks requires domain-specific instruction data, yet
existing generation methods are weakly grounded in source material and
rarely preserve provenance. We present VeritasCarbon, a data-centric
system that transforms 17,721 ESG disclosure documents into 35,009
traceable instruction-response pairs. Its core component, Council of
Domain Experts (CoDE), organizes 11 specialized expert agents in a
four-layer hierarchy, adaptively selects up to K experts per chunk,
coordinates sequential or parallel collaboration, and refines outputs
through MetaExpert-guided feedback. Each retained pair keeps an explicit
mapping to its source chunk, enabling auditing and lineage analysis.
Using the same Qwen2-72B-Instruct backbone for all methods, CoDE on a
2,000-pair evaluation subset substantially outperforms direct prompting,
Self-Instruct, and WizardLM-Evol, each run on 2,000 outputs under the
finalized token-budget protocol: it achieves 4.0x higher ROUGE-L, 4.6x
higher BLEU-4, 10.5x higher Distinct-3, and 14.6x higher domain
relevance than the strongest baselines, while maintaining FactCheck of
0.9478 and near-perfect structural completeness. Ablations show that K =
3 offers the best quality-cost balance, parallel collaboration delivers
the highest overall quality, and feedback remains necessary in the
finalized CoDE protocol. Downstream fine-tuning on CoDE data improves
FactCheck by 2.8x across model architectures (Qwen2.5-7B: 0.185 to
0.510), with five of six models showing consistent gains. VeritasCarbon
demonstrates how provenance-aware, workflow-oriented instruction
generation can support trustworthy domain adaptation in regulated ESG
settings.
PVLDB Artifact Availability:
Permission to make digital or hard copies of all or part of this work
for personal or classroom use is granted without fee provided that
copies are not made or distributed for profit or commercial advantage
and that copies bear this notice and the full citation on the first
page. To copy otherwise, to republish, to post on servers or to
redistribute to lists, requires prior specific permission and/or a fee.
Articles from this volume were invited to present their results at The
53rd International Conference on Very Large Data Bases, August 23--27,
2027, Athens, Greece.
Proceedings of the VLDB Endowment, Vol. 20, No. X
Copyright 2027 VLDB Endowment 2150-8097/27/XX \ldots{} \$10.00.
\end{abstract}
\section{INTRODUCTION}
In recent years, rising global temperatures, widening social inequality,
and high-profile corporate governance failures have intensified pressure
on businesses to operate responsibly. Investors, regulators, and the
public increasingly demand transparent reporting on environmental
impact, social responsibility, and governance practices. As a result,
Environmental, Social, and Governance (ESG) disclosure and carbon
emissions accounting have become critical for organizations facing
growing regulatory scrutiny and societal expectations. Accurate,
transparent, and auditable reporting in these domains is essential for
compliance, risk management, and fostering stakeholder trust. Against
this backdrop, recent advances in large language models (LLMs) have
shown promise in automating complex tasks across specialized domains
such as financial analysis, risk management, and corporate compliance
[21, 53],
[59]. However, deploying LLMs in regulated
environments introduces unique challenges that go beyond those
encountered in open-domain applications.
While prior research demonstrates that LLM performance depends not only
on model scale but also on the quality of prompt data and fine-tuning
methods [3, 39],
[50], these findings primarily reflect open-domain
settings. The instruction tuning paradigm, where models are fine-tuned
on (instruction, output) pairs, has been extensively investigated to
align LLM behavior with user intent [61].
Reinforcement learning from human feedback (RLHF)
[32] and direct preference optimization (DPO)
[36] have further advanced the alignment of LLMs
with human preferences. In such contexts, models are often evaluated on
their ability to generate plausible and diverse responses, and recent
studies suggest that even a small number of high-quality examples can
outperform large-scale noisy datasets [64]. In
contrast, regulated domains like ESG disclosure and carbon emissions
accounting demand strict adherence to methodological rigor, traceable
reasoning, and numerical consistency. For example, organizations must
justify their carbon accounting methods, define clear boundaries, and
ensure that reported figures are auditable and reproducible
[2, 14, 23].
Therefore, although existing LLMs excel at open-ended question
answering, they frequently yield unstable or inconsistent results when
applied to regulatory scenarios where traceability and methodological
consistency are paramount.
A key barrier to reliable LLM deployment in these settings is the lack
of high-quality, instruction-tuning data that enforces source binding
and methodological fidelity. Mainstream instruction data generation
techniques, such as bootstrapping and model distillation
[46, 56], frequently neglect
the necessity for explicit mapping between generated instructions and
their source documents, as well as the requirement for consistent
calculation paths. Representative distillation-based methods, including
Alpaca [41] and Orca [31],
demonstrate that teacher-student pipelines can efficiently produce
large-scale instruction data from proprietary models such as ChatGPT,
yet the generated corpora lack provenance metadata linking outputs to
source evidence. Open-source chat models such as Vicuna
[8] and Baize [57] further
confirm the viability of distillation but inherit the same traceability
limitations. Recent work on automatic data selection for alignment has
shown that filtering instruction data by quality criteria significantly
improves model performance [5]
[28], while self-alignment methods
[26] and unnatural instruction generation
[20] explore alternative paradigms for reducing
human annotation cost. In ESG and carbon accounting, divergent
methodological choices can lead to substantially different outcomes.
Without instruction data that captures these nuances and maintains
source traceability, models risk learning unconstrained or ambiguous
knowledge. Thus, there is a pressing need for instruction data
generation frameworks that ensure both semantic completeness and
traceable source mapping.
From a data management perspective, constructing high-quality
instruction data for ESG applications is a multi-stage transformation
process. Raw ESG disclosure documents are typically lengthy and
heterogeneous, blending narrative text, tabular data, and methodological
explanations. Effective instruction construction requires (1) segmenting
the corpus into semantically complete knowledge units, and (2) ensuring
that each generated instruction can be mapped back to its originating
text fragment for auditing and interpretability. This traceability
requirement is closely related to data provenance research
[7, 18], which emphasizes the
importance of tracking data transformations to support validation and
compliance. The data-centric AI paradigm [52, 60] further highlights the importance of
systematically managing training data, including documentation
[1, 15], quality assessment
[22] [38], and lineage
tracking, is crucial for building reliable AI systems. Model cards
[30] and datasheets [15]
provide complementary documentation frameworks for ensuring transparency
in model development and dataset construction. Without clear definitions
of intermediate data artifacts and their provenance, the instruction
generation process becomes opaque and difficult to audit.
To address these challenges, we introduce VeritasCarbon, a system
designed to generate high-quality, traceable instruction data for ESG
disclosure and carbon emissions accounting within mandatory disclosure
and compliance audit environments. VeritasCarbon formalizes instruction
construction as an optimization problem that balances data quality and
diversity under explicit source constraints. The system processes a raw
ESG corpus through a hierarchical pipeline comprising four stages: (1)
corpus collection and standardization, (2) semantic segmentation and
deduplication, (3) instruction-response generation via expert
scheduling, and (4) quality filtering with source mapping.
Figure~\ref{fig:1} presents an overview of the
VeritasCarbon pipeline. Each stage exposes clear data interfaces,
ensuring controllability and traceability throughout the data generation
lifecycle.
A central component of VeritasCarbon is the Council of Domain Experts
(CoDE) framework, which orchestrates collaborative instruction
generation among multiple expert agents. Drawing on recent advances in
multi-agent LLM systems [6] [19]
[25] [34], CoDE enables adaptive
expert selection and controlled feedback loops, enhancing the stability
and methodological fidelity of generated instruction-response pairs. A
recent survey on LLM-based autonomous agents [48]
identified role specialisation and inter-agent communication as key
enablers of complex task decomposition. Unlike single-agent distillation
approaches, CoDE leverages inter-agent debate and review mechanisms
[13] to improve factual consistency, which is
particularly important in domain-constrained settings. Based on this
framework, we constructed and released the VeritasCarbon-ESG-35K
dataset, comprising 35,000 traceable instruction instances across 11
expert-defined ESG topics. To ensure reproducibility and data security,
we deployed the Qwen2-72B-Instruct [58] 4-bit
quantized model locally, eliminating external API dependencies.
\begin{figure*}[htbp]
\centering
\includegraphics[width=\textwidth]{../results/figures_and_tables/figure1.png}
\caption{The VeritasCarbon pipeline. Raw ESG documents are processed through four stages to produce traceable instruction-response pairs.}
\label{fig:1}
\end{figure*}
We evaluate VeritasCarbon from multiple perspectives. First, we assess
data quality and diversity relative to three baseline methods, employing
both automatic metrics [27, 33, 62] and LLM-as-judge evaluation
[63]. Second, we conduct ablation studies to analyze the impact of
expert selection and feedback mechanisms. Third, we present case studies
to demonstrate the system\textquotesingle s applicability in real-world
ESG scenarios. All experimental protocols and datasets are designed to
facilitate reproducibility.
Our main contributions are as follows:
(1) We introduce a structured workflow with explicit intermediate
artifacts (chunks, expert selections, quality scores) and enforce source
mapping for every generated pair, ensuring full traceability.
(2) We design the CoDE (Council of Domain Experts) framework, which
layer‑selects up to K domain experts, supports sequential and parallel
collaboration, and applies MetaExpert‑guided refinement, which tailored
specifically for ESG disclosure documents.
(3) We build end‑to‑end provenance tracking from the raw corpus through
segmentation, expert routing, collaboration, and filtering, attaching
provenance metadata to each final instruction‑response pair. By
combining these innovations, VeritasCarbon enables reproducible,
auditable instruction data generation for compliance‑critical ESG and
carbon accounting tasks.
(4) We validate CoDE-generated data through downstream fine-tuning on
six diverse model architectures (7B-\/-72B), demonstrating consistent
improvements in source fidelity and factual grounding, with FactCheck
gains of up to 2.76x.
\section{PRELIMINARY}
VeritasCarbon frames ESG instruction-data generation as a constrained
data-transformation problem rather than a purely prompt-level
text-generation task. The system begins with a raw ESG corpus \(C\)
composed of corporate disclosures, regulatory texts, industry reports,
and methodological reference materials. In particular, the disclosure
subset contains heterogeneous corporate ESG documents released by firms
with different business natures, including manufacturing companies,
energy and utility providers, financial institutions, technology firms,
infrastructure operators, logistics companies, and consumer-facing
enterprises. These documents also cover firms operating across different
geographical areas and regulatory environments, including domestic
markets, cross-regional supply chains, and multinational business
contexts. As a result, the corpus reflects diverse reporting
conventions, sector-specific materiality concerns, and
jurisdiction-dependent disclosure requirements.
The ESG disclosure documents used in this study are not uniform
narrative reports. They typically combine several types of information
within the same document: company profiles, business-segment
descriptions, geographic operating footprints, governance structures,
environmental management policies, carbon-emission accounting results,
resource-use indicators, social responsibility programs, supply-chain
practices, risk-management discussions, and references to reporting
standards or regulatory guidelines. For example, an energy company may
emphasize emissions, energy transition, and regulatory compliance; a
manufacturing company may focus on resource efficiency, waste
management, occupational safety, and supply-chain governance; while a
financial institution may disclose sustainable finance, climate-risk
exposure, green investment policies, and governance procedures. This
heterogeneity makes ESG disclosures a challenging but valuable source
for instruction-data generation.
After parsing, cleaning, segmentation, and deduplication, the corpus is
converted into a set of semantically coherent chunks \{\(c_{1}\),
\(c_{2}\), ..., \(c_{n}\)\}. Each chunk is not treated as an isolated
text span only; it is associated with contextual metadata such as
document type, source organization, business sector, reporting year,
geographical scope, and provenance identifier when such information is
available. These contextual attributes help the system distinguish
whether a chunk describes a company-level sustainability policy, a
region-specific regulatory requirement, a sector-specific
carbon-accounting practice, or a general methodological reference. They
also support more precise expert selection and reduce the risk of
generating responses that ignore the business or geographical context of
the disclosure.
For each chunk \(c_{i}\), the system dynamically selects a subset of
domain expert agents from an expert pool E and generates one or more
candidate instruction-response pairs. The selected experts interpret
both the local textual evidence and the available contextual metadata.
This is important because many ESG claims are context-dependent: a
disclosure about renewable-energy procurement, for instance, may have
different implications for a utility provider, a manufacturing firm, or
a financial institution; similarly, carbon-accounting statements may
differ depending on whether they refer to direct operations, purchased
electricity, value-chain emissions, or regional regulatory obligations.
The final dataset \(D\) is stored as tuples (\(q_{j}\), \(a_{j}\),
\(s_{j}\)), where \(q_{j}\) is the instruction, \(a_{j}\) is the
response, and \(s_{j}\) denotes the source mapping back to one or more
chunks. In the extended representation, each tuple can also include
metadata fields such as document category, sector, reporting period,
geographical scope, expert path, and quality score. This representation
makes provenance a first-class data object rather than a post hoc
annotation. It also allows users to inspect not only what answer was
generated, but also from which disclosure context the answer was derived
and which expert perspectives contributed to its construction.
Three design constraints govern the entire pipeline. First, source
fidelity: every response must remain grounded in the originating text
fragment and must preserve the factual boundaries of the source. Second,
diversity: the generated instructions should cover multiple ESG themes,
task forms, and reasoning styles rather than collapse into repetitive
question templates. Third, quality: only outputs that exceed a minimum
threshold are retained. Formally, the optimization target is to maximize
expected dataset quality \(E\lbrack Q(D)\rbrack\) while satisfying
source-mapping consistency and topical diversity constraints. This
formalization is particularly important in ESG and carbon-accounting
settings because the same disclosure often interweaves narrative
explanation, numerical evidence, standards reference, and implicit
assumptions.
\begin{table}[htbp]
\centering
\small
\caption{Core symbols and execution variables.}
\label{tab:1}
\begin{tabular}{@{}>{\raggedright\arraybackslash}p{0.12\columnwidth}>{\raggedright\arraybackslash}p{0.25\columnwidth}>{\raggedright\arraybackslash}p{0.45\columnwidth}@{}}
\toprule
\textbf{Symbol} & \textbf{Description} & \textbf{Meaning} \\
\midrule
\(c_i\) & Source chunk & A semantic segment of an ESG document \\
\(E\) & Expert pool & Set of 11 specialized expert agents \\
\(K\) & Max experts per chunk & Upper bound on active experts (default 3) \\
\(R\) & Max feedback rounds & Upper bound on MetaExpert iterations (default 2) \\
\(Q(o,c)\) & Quality function & Composite score of output \(o\) on chunk \(c\) \\
\(\tau\) & Release threshold & Minimum quality score for dataset inclusion \\
\(D\) & Final dataset & Traceable instruction-response pairs with metadata \\
\bottomrule
\end{tabular}
\end{table}
The published release contains 35,009 traceable instruction-response
pairs generated from 20,000 source chunks and 17,721 ESG disclosure
documents. Table~\ref{tab:2} summarizes the key
statistics of the VeritasCarbon-ESG-35K dataset. The corpus is arranged
into four layers (foundational methodology texts, public corporate
reports, policy documents, and industry analysis) to support both broad
domain coverage and adaptive expert routing. The dataset also records a
45.1\% multi‑expert collaboration ratio and a 100\% MetaExpert
participation ratio in the final released set.
\begin{table}[htbp]
\centering
\small
\caption{Corpus scope.}
\label{tab:2}
\begin{tabular}{@{}>{\raggedright\arraybackslash}p{0.45\columnwidth}>{\raggedright\arraybackslash}p{0.45\columnwidth}@{}}
\toprule
\textbf{Dataset Statistic} & \textbf{Value} \\
\midrule
Total instruction-response pairs & 35,009 \\
Source chunks & 20,000 \\
Source documents & 17,721 \\
Corpus layers & 4 \\
Expert types & 11 \\
Multi-expert ratio & 45.1\% \\
Average quality score & 0.667 ($\sigma$= 0.103) \\
Average instruction length & 106.5 characters \\
Average response length & 380.4 characters \\
\bottomrule
\end{tabular}
\end{table}
\section{THE VERITASCARBON FRAMEWORK}
At the system level, VeritasCarbon is implemented as a modular
four-stage pipeline with explicit interfaces between collection, corpus
engineering, expert-guided generation, and release filtering. This
modularization is central to the paper's data-management contribution:
every stage emits an intermediate artifact that can be audited, re-run,
or replaced without collapsing the entire workflow. In contrast to
generic instruction bootstrapping methods that optimize only the final
generated text, VeritasCarbon optimizes the full path from source
document to released dataset record.
\textbf{3.1 Stage-level architecture}
Stage 1, corpus collection and standardization, converts heterogeneous
ESG documents into a unified processing representation. Stage 2,
semantic segmentation and deduplication, transforms long semi-structured
documents into semantically complete chunks suitable for expert routing.
Stage 3, expert-guided generation, activates the CoDE engine to select
experts, compose interaction patterns, produce candidate QA pairs, and
refine them over bounded feedback rounds. Stage 4, quality filtering and
source mapping, applies consistency checks and persists only those
outputs that satisfy release criteria. The pipeline image also
highlights an important system property: the dataset is not created by a
single monolithic generator. It is created by a controlled workflow in
which routing, selection, generation, scoring, and packaging are
distinct operations.
This architecture supports controllability and reproducibility. Because
the pipeline is hierarchical and interface-driven, it can be debugged as
a workflow rather than treated as an opaque prompt chain. That design
choice matters in ESG settings, where downstream users need not only
strong responses but also a defensible trail back to source evidence and
transformation history.
3.2 Internal CoDE organization
The core engine of VeritasCarbon is CoDE (Council of Domain Experts).
Figure~\ref{fig:2} illustrates its internal
architecture, which comprises seven interconnected modules.
The workflow begins with Module\,1, which receives an input ESG chunk
along with its metadata, source type, provenance ID, and task objective.
Module\,2 is a layered expert pool containing 11 agents organized into
three tiers: general understanding (e.g., Disclosure, Policy), domain
reasoning (Carbon Accounting, Risk, Governance), and quality assurance
(Materiality, Assurance, Data Quality). Module\,3, the MetaExpert
orchestrator, parses the chunk intent, selects relevant expert layers,
composes a team of up to \(K\) experts, and schedules the collaboration
mode (sequential, parallel, or hybrid). Module\,4 is the collaboration
workspace, where selected experts generate draft questions and answers,
critique each other\textquotesingle s outputs, and align responses with
source evidence. Module\,5 implements a feedback refinement loop,
applying structural, factual, and domain‑relevance checks; the loop
iterates until quality is accepted or the maximum number of rounds R is
reached. Module\,6 performs verification and traceability packaging,
scoring domain relevance, factual grounding, structural completeness,
and non‑redundancy, then attaching source‑span links, expert‑path logs,
and provenance metadata. Finally, Module\,7 outputs a traceable QA pair
containing the instruction, response, source evidence, provenance
metadata, expert path, and quality score.
\begin{figure*}[htbp]
\centering
\includegraphics[width=\textwidth]{../results/figures_and_tables/figure2.png}
\caption{CoDE internal architecture. Layered expert selection, MetaExpert orchestration, collaboration modes, and feedback refinement.}
\label{fig:2}
\end{figure*}
\textbf{3.3 Provenance as a framework invariant}
A distinguishing framework invariant is that provenance is preserved
throughout the pipeline. Every released pair remains bound to its
originating chunk or chunk set, and the workflow logs intermediate
expert outputs as well as feedback iterations. In practical terms, this
means the framework can answer not only \textquotesingle what
instruction-response pair was released?\textquotesingle{} but also
\textquotesingle where did it come from?\textquotesingle,
\textquotesingle which experts participated?\textquotesingle,
\textquotesingle which quality gate admitted it?\textquotesingle, and
\textquotesingle what intermediate transformations occurred before
release?\textquotesingle. This is the point at which VeritasCarbon moves
from being a prompt recipe to being a data-management system.
\section{METHODS}
This section rewrites the core technical procedure in
implementation-oriented language. The intention is to make the
bottom-layer logic easier to follow for readers who need to reproduce,
adapt, or extend the system.
\textbf{4.1 Corpus collection, parsing, and normalization}
The corpus processing layer ingests ESG reports in multiple file formats
and converts them into a unified textual representation. The published
paper states that the corpus spans four layers: 97 domain textbooks and
reference materials, 17,425 CSR reports, 92 policy documents, and 107
industry research reports. Format-specific extractors preserve
structural signals such as headings, tables, and paragraph boundaries
during parsing. This detail is important because later segmentation and
source mapping assume that structural boundaries still carry semantic
value.
The accompanying preprocessing notebook extends that description with
implementation details. The preprocessing routine supports PDF, DOCX,
and TXT inputs; performs text cleaning to remove headers, footers, page
numbers, and structural noise; repairs garbled encodings; converts
Traditional Chinese to Simplified Chinese where needed; executes line-,
paragraph-, and chunk-level deduplication; applies quality gates such as
valid-character checks; and can be checkpointed and resumed. The
notebook also records a cleaned chunk file, an optional audit log, and a
checkpoint file to support reproducible reruns. These engineering
details strengthen the reproducibility story because they make
preprocessing a persistent data product rather than an undocumented
one-time script.
\textbf{4.2 Semantic segmentation and deduplication}
After parsing, long ESG documents are split into semantically complete
chunks using a sliding-window strategy with overlap. Each chunk is
constrained to roughly 200-800 characters and that split points are
aligned with paragraph and section boundaries whenever possible. This
choice balances two competing requirements: chunks must be short enough
for efficient LLM processing, but long enough to preserve the local
context needed for source-grounded generation.
Near-duplicate chunks are then identified with character n-gram
similarity. Chunks whose Jaccard similarity exceeds 0.85 are merged, and
the more complete version is retained. The goal is not merely storage
reduction; it is to prevent repeated or nearly repeated content from
biasing generation and evaluation. By deduplicating early, the system
reduces the chance that one disclosure pattern dominates the instruction
distribution simply because it is repeated across pages or documents.
\textbf{4.3 Layered expert selection}
Given a chunk ci, VeritasCarbon computes a feature vector fi describing
its content characteristics, such as whether the chunk contains
numerical evidence, temporal references, standards language,
comparisons, or structured entity-relation patterns. Expert activation
is then performed layer by layer. The base layer includes QA, Summary,
Extraction, Classification, and Analysis experts, and at least one
base-layer expert is always selected. The analysis layer contains
Temporal Analysis, Benchmark Comparison, and Greenwashing Detection
experts, which are activated when the chunk exhibits temporal,
comparative, or potentially misleading disclosure patterns. The
verification layer contains Consistency Verification and Standard
Alignment experts for chunks that require numerical cross-checking or
reporting-standard alignment. The graph layer contains the Knowledge
Graph expert, which is activated when relation-rich entity structure is
detected.
The routing policy accumulates experts across layers and truncates the
active set to a maximum of \(K\) experts, with \(K\) = 3 used as the
default operational setting. If the selection procedure returns too few
experts, the controller adds additional base-layer experts to guarantee
minimum coverage. Operationally, this makes CoDE a sparse-activation
system: straightforward chunks remain inexpensive, while complex chunks
invoke deeper specialization only when their observed features justify
the extra computational cost.
\textbf{4.4 Collaboration modes and candidate construction}
Once the active expert set has been fixed, VeritasCarbon supports two
collaboration modes. In sequential collaboration, experts process the
chunk in an ordered chain, and each expert receives the output of the
previous expert as added context. This mode encourages progressive
refinement, because later experts can correct, sharpen, or specialize an
earlier draft. In parallel collaboration, all selected experts process
the same chunk independently. The controller then compares the generated
candidates with the quality function \(Q\ (o,c)\) and retains the best
candidate. This mode encourages diversity by allowing competing
interpretations to be explored simultaneously before the system commits
to one response.
Both models are valid from a systems perspective. Sequential
collaboration is closer to a staged refinement workflow, whereas
parallel collaboration resembles a best-of-K execution strategy. Results
show that both are beneficial relative to no collaboration, with
parallel mode reaching the highest mean quality and sequential mode
retaining a slight BLEU-4 advantage.
\textbf{4.5 MetaExpert orchestration and feedback refinement}
The MetaExpert is the coordinating component that converts expert
outputs into releasable candidates. It operates in two stages. First, it
aggregates and ranks expert outputs, then selects the top-k candidate
drafts according to the quality scoring function. Second, it runs an
iterative feedback loop over those candidates. In each round, the
MetaExpert emits targeted improvement signals, such as requests for
tighter source anchoring, better instruction formulation, or improved
alignment with ESG topic intent. A quality gate is applied after each
iteration, and the loop terminates either when the candidate converges
or when the maximum number of rounds R is reached.
When R = 0 still produces valid outputs, but with lower quality than R =
2. That result matters methodologically because it shows that the
architecture is not dependent on feedback for basic generation, yet
still benefits substantially from bounded iterative
refinement.\(R\ \)\(R\)
\textbf{4.6 Quality validation and traceability packaging}
After generation, candidates move into a validation-and-packaging phase.
The quality validator scores domain relevance, source fidelity, lexical
diversity, factual grounding, and structural completeness. Only
candidates that exceed the release threshold $\tau$ are preserved. The
traceability binder then attaches source-chunk identifiers, expert-path
metadata, and related provenance records so that each dataset item can
be traced back to its originating corpus fragment.
The traceability binder formalizes provenance as a structured schema
rather than an ad-hoc annotation. Each released record carries a
provenance record P = (source\_chunk\_id, expert\_path,
collaboration\_mode, feedback\_rounds, quality\_score, timestamp). The
source\_chunk\_id is a composite key encoding the corpus layer, document
identifier, and chunk index, enabling reverse lookup from any dataset
record to its exact textual origin in the source corpus. The
expert\_path records the ordered list of expert agents that contributed
to the generation, including their selection scores and whether each
expert\textquotesingle s output was retained in the final response or
discarded during the voting/critique phase. This structured approach
means that provenance queries-\/-\/-such as "find all QA pairs where the
Greenwashing Detection expert was activated and the quality score
exceeds 0.7"-\/-\/-are directly answerable from the dataset without
recomputing the generation pipeline. Schema details and query examples
are provided in Appendix A.7.
The quality validator applies a composite scoring function Q(o, c) that
returns a vector of four sub-scores: correctness (weight 0.4),
completeness (0.3), relevance (0.2), and structure (0.1). We note that
this same scoring rubric is used for both generation-time feedback
(Section 4.5) and release-time filtering (this section). To avoid
circularity in evaluation, the intrinsic metrics reported in Section 5.2
and the downstream fine-tuning evaluation in Section 5.6 are computed
independently of Q(o, c). The MetaExpert\textquotesingle s quality score
is retained in the dataset as a generation-time diagnostic but is not
used as an evaluation metric in the paper\textquotesingle s experimental
sections.
\textbf{4.7 Reproducibility assets and practical reruns}
The paper and the associated files emphasize reproducibility at three
levels. At the data level, preprocessing outputs are materialized as
cleaned chunk files and audit logs. At the model level, all compared
methods use the same Qwen2-72B-Instruct model so that performance
differences can be attributed to method design rather than model choice.
At the workflow level, the system can be re-executed from intermediate
artifacts, which is essential for debugging quality regressions,
re-running only selected corpus layers, or auditing changes in
expert-selection logic. Together, these choices make VeritasCarbon more
reproducible than prompt-only pipelines whose intermediate
transformations are discarded.
\section{EVALUATION AND PERFORMANCE}
The evaluation is organized around four questions: whether CoDE improves
intrinsic data quality relative to strong baselines, which architectural
choices drive performance, whether the system scales with corpus size,
and what the released dataset looks like as a data product. All methods
use the same underlying model (Qwen2-72B-Instruct) and are measured on
aligned samples to ensure fair comparison.
5.1 Experimental setup
The released VeritasCarbon-ESG-35K dataset is constructed from 17,721
source documents and 20,000 semantically complete chunks. CoDE generates
35,009 instruction-response pairs, with all 11 expert types represented
and 45.1\% of released pairs involving multi-expert collaboration.
Baseline comparisons are conducted on n = 2,000 pairs per method using
direct prompting, Self-Instruct, and WizardLM-Evol. Ablation experiments
are performed on n = 500 samples per condition across four design
dimensions: expert count, collaboration mode, feedback rounds, and
knowledge injection.
Seven intrinsic metrics are used in the paper. ROUGE-L measures source
fidelity through longest-common-subsequence overlaps with the source
chunk. BLEU-4 measures precision of generated wording. Distinct-2 and
Distinct-3 measure lexical diversity. Domain Relevance captures coverage
of ESG-specific vocabulary. FactCheck measures how often numerical
values in the response appear in the source chunk. Structural
Completeness checks that instructions and responses are non-trivial in
length. Together, these metrics emphasize linguistic quality, grounding,
domain fit, and release readiness.
5.2 Main comparative results
The main intrinsic evaluation demonstrates a wide margin between CoDE
and all three baselines (Figure 3). The largest gaps appear in source
fidelity, domain relevance, and generation diversity. CoDE achieves
ROUGE‑L = 0.3380 and BLEU‑4 = 0.1932, compared with 0.0838 and 0.0424
for direct prompting. It also reaches domain relevance = 0.3637, far
above the best baseline at 0.0249. These gains confirm that
domain‑specialized multi‑expert generation produces outputs more
strongly anchored to ESG source material than generic
instruction‑generation pipelines.
\begin{figure}[htbp]
\centering
\includegraphics[width=\columnwidth]{../results/figures_and_tables/fig5_coe_vs_baselines.png}
\caption{Comparison of CoDE against three baselines.}
\label{fig:4}
\end{figure}
The baseline breakdown is also informative. Self-Instruct maintains
relatively strong FactCheck but very low ROUGE-L, suggesting safer yet
weakly grounded responses. WizardLM-Evol retains some domain vocabulary
but performs poorly on factual grounding. Direct Prompting is the
strongest baseline overall, yet still falls well short of CoDE on every
quality-oriented metric except structural completeness, where all
methods are near saturation.
5.3 Ablation study
The ablation study isolates the contribution of each main design
dimension (Table 3). With respect to expert count, moving from K = 1 to
K = 2 produces the largest jump in mean quality, while gains beyond K =
2 become marginal. Therefore, we interpret K = 3 as a practical
compromise between output quality and cost. For collaboration mode, both
sequential and parallel collaboration outperform the no-collaboration
baseline, confirming that expert interaction is intrinsically beneficial
rather than decorative. Parallel collaboration attains the highest
reported mean quality (0.647). For feedback rounds: when feedback is
disabled (R = 0), the system still produces valid outputs, but mean
quality rises from 0.626 at R = 0 to 0.649 at R = 2. Thus, feedback acts
as a quality amplifier rather than an architectural prerequisite.
Knowledge injection shows a smaller but positive effect, indicating that
prompt‑specialized experts already encode meaningful domain knowledge,
yet explicit injection remains useful for alignment with ESG terminology
and standards.\(K\)\(K\)\(K\)\(K\)\(R\)\(R\)\(R\)
5.4 Scalability evaluation
The scalability study varies the input size from 500 to 20,000 chunks.
Figure~\ref{fig:5} plots throughput as a function of
corpus size. Wall-clock time grows approximately linearly, while
throughput remains stable between 23 and 26 QA pairs per minute.
Processing 20,000 chunks takes about 12.8 hours, and token consumption
scales linearly at roughly 800 tokens per chunk. These results confirm
that CoDE is suitable for large-scale data generation without
architectural modifications. Table 4 reports detailed resource
consumption (generation time, throughput, API calls, and tokens) across
corpus sizes.
\begin{figure}[htbp]
\centering
\includegraphics[width=\columnwidth]{../results/scalability/fig_scalability_corpus_size.png}
\caption{Generation time scales linearly with corpus size. (a) Scalability of CoDE; (b) throughput stability.}
\label{fig:5}
\end{figure}
Table 3 Ablation study results across four dimensions ($n$ = 500 per
condition). Best results per dimension are in bold.
5.5 Dataset analysis and release characteristics
Figure~\ref{fig:6} shows the distribution of quality
scores across the 35,009 generated QA pairs. The distribution centers
around 0.667 with \(\sigma\) = 0.103, and 89.7\% of pairs exceed the
minimum threshold \(\tau\) = 0.5.
The expert‑usage profile is highly skewed in a meaningful way.
Figure~\ref{fig:7} presents the frequency of each expert
agent's invocation. The Summary expert dominates overall use, followed
by Consistency Verification and Classification. This pattern supports
the system's sparse‑activation logic: base‑layer experts handle the
majority of chunks, while upper‑layer specialists such as Greenwashing
Detection, Temporal Analysis, Benchmark Comparison, and Knowledge Graph
are invoked only when the content indicates a need for them. In other
words, the expert distribution is not balanced by design; it is adaptive
by design.
\begin{figure}[htbp]
\centering
\includegraphics[width=\columnwidth]{../results/figures_and_tables/fig2_quality_distribution.png}
\caption{Distribution of quality scores across 35,009 generated QA pairs.}
\label{fig:6}
\end{figure}
\begin{figure}[htbp]
\centering
\includegraphics[width=\columnwidth]{../results/figures_and_tables/fig1_expert_usage.png}
\caption{Distribution of expert agent invocations across 35,009 QA pairs. Summary Expert dominates base-layer usage, while Consistency Verification Expert is the most frequently activated upper-layer agent.}
\label{fig:7}
\end{figure}
\textbf{5.6 Downstream Fine-tuning Evaluation}
While intrinsic metrics (Section 5.2-\/-5.3) measure the quality of
generated instruction data, the ultimate test is whether this data
improves downstream ESG task performance when used for supervised
fine-tuning. We therefore conduct a cross-model fine-tuning evaluation
to validate that CoDE-generated data translates to measurable model
improvements.
Experimental Protocol. We fine-tune six instruction-tuned models
spanning three architecture families and three parameter scales on
31,315 CoDE training pairs (filtered subset of VeritasCarbon-ESG-35K).
All models are fine-tuned with QLoRA (r=16, alpha=32, 4-bit NF4
quantization) for 3 epochs using the same Alpaca-style prompt template.
We evaluate on a held-out ESG QA test set (n=200) using the same seven
intrinsic metrics defined in Section 5.1. For each model, we compare the
fine-tuned variant against its zero-shot counterpart. Statistical
significance is assessed via bootstrap with 95\% confidence intervals.
Baseline Models. We select models spanning three regulatory-relevant
dimensions: (i) model scale (7B to 14B to 72B within the Qwen2.5
family), (ii) architecture diversity (dense Qwen/InternLM vs. MoE
DeepSeek-V2-Lite), and (iii) training provenance (Chinese-native
GLM-4-9B vs. bilingual Qwen/InternLM/DeepSeek). All models are
instruction-tuned chat variants, ensuring that any improvement from
fine-tuning is attributable to domain knowledge acquisition rather than
instruction-following capability.
Fine-tuning Results. Table 5 reports the main results. Across five of
six models, fine-tuning on CoDE data produces consistent and
statistically significant improvements. The key findings are:
(1) Source fidelity improves across scales. Qwen2-72B fine-tuned
achieves ROUGE-L = 0.2441 (vs. 0.1132 zero-shot, a 2.2x improvement),
BERTScore-F1 = 0.7264, and the highest BLEU-4 among all models (0.1307).
The 7B variant similarly gains +23.7\% ROUGE-L (0.1945 to 0.2406, p
\textless{} 0.001).
(2) Factual grounding improves dramatically. The most striking gain is
FactCheck, which measures numerical overlap between generated responses
and source evidence. Fine-tuned Qwen2.5-7B achieves FactCheck = 0.510
(vs. 0.185 zero-shot, a 2.76x improvement). InternLM2.5-7B shows similar
magnitude (0.230 to 0.630, 2.74x). These gains confirm that CoDE data
successfully trains models to anchor numerical claims in source
evidence-\/-\/-a critical capability for ESG reporting.
(3) Domain relevance is near-ceiling. All fine-tuned models exceed 0.90
domain relevance (Qwen2.5-7B-FT: 0.997), confirming that CoDE data
effectively encodes ESG domain vocabulary and concepts.
(4) Not all architectures benefit equally. GLM-4-9B exhibits negative
transfer: fine-tuned ROUGE-L drops from 0.1953 to 0.1582, and
BERTScore-F1 degrades from 0.6981 to 0.5216. We attribute this to
tokenizer mismatch between CoDE data (generated with
Qwen2-72B\textquotesingle s predominantly Chinese-centric tokenizer) and
GLM-4-9B\textquotesingle s bilingual tokenizer, resulting in suboptimal
sequence packing when processing Chinese ESG terminology. This finding
underscores the importance of tokenizer compatibility in domain-specific
instruction tuning.
LLM-as-Judge Validation. To complement automatic metrics with
qualitative assessment, we conduct an LLM-as-judge evaluation [63]
on n=100 generated responses per model. A strong evaluator model (GPT-4)
rates each response on a 5-point Likert scale along three dimensions:
(i) accuracy-\/-\/-whether claims are factually correct and grounded in
the source, (ii) domain relevance-\/-\/-whether ESG terminology and
concepts are used appropriately, and (iii) fluency-\/-\/-whether the
response is well-structured and coherent. The overall score is the
unweighted mean of the three dimensions.
Qwen2-72B fine-tuned achieves the highest overall judge score of 4.46/5
(Accuracy: 4.32, Domain Relevance: 4.55, Fluency: 4.52), followed by
Qwen2.5-7B fine-tuned at 4.38/5. Notably, LLM-as-Judge rankings are
broadly consistent with automatic metric rankings, providing convergent
validation. The zero-shot Qwen2-72B already scores 4.35, indicating
strong innate ESG capability; fine-tuning on CoDE data provides an
additional quality margin even on top of this strong baseline. GLM-4-9B
fine-tuned scores only 1.96/5 overall, confirming the catastrophic
degradation observed in automatic metrics and reinforcing the tokenizer
compatibility hypothesis.
\emph{\textbf{Table 5: Fine-tuning evaluation on ESG domain test set
(n=200), with LLM-as-Judge scores (n=100 per model).}}
Taken together, these results demonstrate that CoDE-generated
instruction data is not merely intrinsically higher-quality (Section
5.2) but reliably transfers to measurable downstream gains across
diverse model architectures and scales. The tokenizer compatibility
failure with GLM-4-9B also highlights a practical lesson for
practitioners deploying domain-specific instruction data: model-data
tokenizer alignment should be explicitly verified before fine-tuning.
\textbf{5.7 LLM-as-Judge Quality Assessment}
To complement automatic metrics, we conduct an LLM-as-judge evaluation
[63] on n=100 generated responses per fine-tuned model. GPT-4 rates
each response on a 5-point Likert scale across three dimensions: (i)
accuracy-\/-\/-whether claims are correct and source-grounded, (ii)
domain relevance-\/-\/-whether ESG terminology is used appropriately,
and (iii) fluency-\/-\/-structure and coherence.
\emph{\textbf{Table 6: LLM-as-Judge evaluation (n=100 per model, 5-point
Likert scale).}}
Key findings: (1) Qwen2-72B-FT achieves the highest overall score
(4.46/5, Accuracy: 4.32, Domain: 4.55, Fluency: 4.52), confirming the
automatic metric rankings. (2) Fine-tuning consistently improves domain
relevance across architectures (Qwen2.5-7B: 4.18 to 4.42;
InternLM2.5-7B: 3.70 to 4.17). (3) GLM-4-9B-FT drops to 1.96/5 overall,
corroborating the catastrophic negative transfer observed in automatic
metrics. (4) Judge scores show convergent validity with automatic
metrics-\/-\/-the evaluation protocol is robust.
This LLM-as-judge evaluation serves as a scalable proxy for human
evaluation, which is deferred to a companion study focusing on
domain-specific ESG assessment with expert annotators.
\textbf{5.8 Limitations of Intrinsic Metrics}
We note an important methodological nuance regarding the intrinsic
evaluation in Section 5.2. ROUGE-L and BLEU-4, as employed here, measure
lexical overlap between generated responses and source chunks rather
than against human-written reference answers. This usage reflects our
evaluation objective-\/-\/-source fidelity-\/-\/-but diverges from
standard practice where these metrics compare system outputs against
gold-standard references. Under this definition, CoDE exhibits an
inherent advantage because it explicitly optimizes for source anchoring
through expert routing and MetaExpert refinement, whereas the baseline
methods are not designed with this constraint.
This asymmetry motivates two design choices in our evaluation protocol.
First, we report FactCheck (numerical overlap with source) as a
complementary grounding metric that is less sensitive to lexical overlap
than ROUGE-L or BLEU-4. Second, and more importantly, the downstream
fine-tuning evaluation (Section 5.6) provides an orthogonal validation
that does not depend on source-chunk comparisons: fine-tuned models are
evaluated on their ability to generate correct ESG responses, and the
consistent improvement over zero-shot baselines confirms that CoDE data
quality is not an artifact of the source-fidelity metric definition. We
encourage future work to develop standardized ESG-domain evaluation
benchmarks with human-annotated reference answers.
\textbf{5.9 Performance interpretation, limitations, and reuse guidance}
Taken together, the results show that CoDE improves the aspects of
generation that matter most for ESG data management: stronger source
fidelity, markedly higher ESG-domain relevance, higher diversity, and
high factual grounding. The system also scales without architectural
modification, which makes it more valuable as a reusable
data-engineering asset. At the same time, there are several limitations.
Intrinsic metrics do not completely determine downstream task utility.
Expert prompts may introduce style or coverage bias. And although the
system scales to 20,000 chunks, additional optimization would be
necessary for real-time or streaming deployment scenarios.
For reuse, the most defensible practice is to preserve the pipeline's
intermediate artifacts: cleaned chunks, routing decisions, expert
outputs, quality scores, and provenance metadata. Doing so allows future
researchers to extend the corpus, swap in new experts, alter the
selection thresholds, or re-tune the quality gate without losing
auditability. This is especially important if future versions add
cross-document synthesis, multilingual processing, or learned quality
estimators.
\section{RELATED WORK}
\subsection{Instruction Data Generation}
High-quality instruction data is foundational for the effective training
and tuning of LLMs, particularly in specialized domains such as ESG and
carbon accounting. A comprehensive survey of instruction tuning methods
[61] identifies the (instruction, output)
supervised fine-tuning paradigm as a critical bridge between the
next-word prediction objective of LLMs and the goal of following user
instructions. Early work based on GPT-3 [3]
demonstrated the value of large-scale pre-training, and subsequent
studies showed that instruction tuning substantially improves zero-shot
and cross-task generalization [39]
[50]. Reinforcement learning from human feedback
(RLHF) [32] and direct preference optimization
(DPO) [36] further align LLMs with human
preferences without explicit reward models.
To reduce manual annotation costs, bootstrapping methods have emerged.
Self-Instruct [46] iteratively expands seed tasks
using model-generated instructions, and WizardLM
[56] increases instruction complexity through
evolutionary prompting. In the distillation paradigm, Alpaca
[41] fine-tunes open-source models (e.g., LLaMA
[42], Llama 2 [43]) on 52K
examples distilled from proprietary models. Orca
[31] learns from rich explanation traces of GPT-4
rather than shallow imitation, while Zephyr [44]
explores distilled preference optimization. Vicuna
[8] and Baize [57]
fine-tune on user-shared or self-generated conversations. Unnatural
instructions [20] and instruction backtranslation
[26] further reduce human effort. A systematic
comparison of 12 instruction datasets [47] shows
that no single source dominates across all evaluations.
Recent studies on data quality challenge the assumption that larger
datasets are always better. LIMA [64] shows that
1,000 carefully curated examples can match models trained on orders of
magnitude more data. AlpaGasus [5] filters
low-quality instructions to improve performance, and Liu et al.
[28] provide a comprehensive study of automatic
data selection for instruction tuning.
Therefore, existing methods focus on open-domain linguistic diversity,
task coverage, and reasoning complexity. Instruction generation is
typically treated as an end-to-end process, producing outputs directly
from prompts without explicit, structured mappings to source documents.
Consequently, intermediate artifacts and source-level constraints are
not modeled within the generation workflow. The absence of provenance
metadata and traceability guarantees limits their applicability in
regulated domains like ESG, where auditability is critical.
\subsection{Domain-Specific Large Language Models}
Adapting LLMs to specialized domains has received growing attention,
particularly in finance and sustainability [60].
BloombergGPT [53] is trained on large-scale
proprietary financial corpora, while FinGPT [59]
provides open-source data pipelines and lightweight adaptation (e.g.,
LoRA). PIXIU [55] offers a comprehensive
framework including a financial instruction dataset (136K samples), a
fine-tuned model (FinMA), and a benchmark covering five tasks. FinBERT
[21] extracts structured information from
financial text.
For climate and sustainability, ClimateBERT [49]
fine-tunes pre-trained models for climate-related text understanding and
has been applied to analyze corporate climate risk disclosures
[2]. Kölbel et al. [23] use
BERT-based models to study how regulatory disclosure of climate risks
affects credit default swap term structures. Friederich et al.
[14] automate identification of climate risk
disclosures in annual reports, while ClimaText
[45] provides a dataset for climate change topic
detection. Luccioni et al. [29] estimate the
carbon footprint of training BLOOM, linking LLM research to
environmental sustainability.
\subsection{Data Provenance, Quality, and Traceability}
The concept of data provenance, which involves the tracking of the
origin, transformation, and lineage of data, has a long history in the
field of database research. Seminal work by Buneman et al.
[4] characterizes the "why" and "where" dimensions
of data provenance. Green et al. [17] formalize
provenance using semiring annotations. Davidson and Freire
[11] survey provenance challenges in scientific
workflows, while Cui et al. [10] address lineage
tracing in data warehousing environments. Cheney et al.
[7] provide a comprehensive treatment of
provenance in databases covering foundational theory and practical
considerations. Moreover, Herschel et al. [18]
offer a modern survey categorizing provenance by purpose,
representation, and source.
Data quality research complements provenance. Rahm and Do
[37] provide an early taxonomy of data cleaning;
Chu et al. [9] survey emerging challenges. Ilyas
and Chu [22] give a textbook treatment, and
HoloClean [38] uses probabilistic inference for
holistic data repair. In data-centric AI, Whang et al.
[52] argue that data quality management must be
integrated throughout the ML pipeline, not treated as a preprocessing
step.
For machine learning pipelines, Polyzotis et al.
[35] identify data management challenges in
production ML systems, and Doan and Halevy [12]
highlight semantic integration issues. Tools for automatically tracking
ML experiment metadata and provenance have been developed
[40]. Grafberger et al.
[16] address data distribution debugging. At the
dataset level, Datasheets for Datasets [15], Data
Statements [1], and Model Cards
[30] provide documentation frameworks for
transparency.
Our work draws on these provenance and quality foundations but extends
them to a novel setting: tracking lineage from source ESG documents
through intermediate knowledge units to final instruction-response
pairs, ensuring each generated instance is auditable and traceable to
its documentary origin.
\subsection{Multi-Agent Collaboration for Data Generation}
Recent LLM-based multi-agent systems enable multiple model instances
with specialized roles to solve complex tasks. A survey by Wang et al.
[48] categorizes LLM-based autonomous agents
along dimensions of profile, memory, planning, and action, providing a
structured taxonomy for understanding agent architectures. AutoGen
[54] provides a general-purpose framework for
multi-agent conversation, enabling flexible orchestration of LLM agents.
CAMEL [25] introduces a role-playing
communicative agent framework that uses inception prompting to guide
autonomous cooperation between agents, demonstrating the potential for
scalable multi-agent data generation. MetaGPT
[19] assigns specialized roles (product manager,
architect, engineer) within a software development workflow. Park et al.
[34] demonstrate ``generative agents'' that
simulate believable human behavior. AgentVerse [6]
facilitates collaboration and explores emergent group behaviors. Du et
al. [13] show that multi-agent debate, where
agents argue and critique each other, improves factuality and reasoning.
In parallel, prompting strategies such as Chain-of-Thought
[51] decompose reasoning into intermediate steps
to enhance output quality at the inference level. Retrieval-augmented
generation [24] grounds LLM outputs in retrieved
evidence, improving factual accuracy. For evaluation, LLM-as-judge
approaches [63] offer scalable alternatives to
human evaluation, employing strong models to assess the quality of
generated outputs.
These multi-agent techniques primarily focus on improving reasoning or
task-solving capability. The integration of multi-agent collaboration
with structured workflow management, intermediate artifact control, and
explicit provenance tracking within generative instruction construction
remains underexplored. Bridging multi-agent LLM systems with data
management practices is crucial for domains that require traceability,
auditability, and compliance.
\subsection{Summary and Positioning}
To summarize, existing approaches suffer from four key limitations: (1)
absence of provenance metadata linking instructions to source documents;
(2) lack of domain-specific instruction datasets built with regulatory
compliance in mind; (3) provenance and quality techniques not embedded
into LLM data generation pipelines; and (4) multi-agent systems
disconnected from workflow-level artifact management.
\section{CONCLUSION AND FUTURE WORK}
We presented VeritasCarbon, a provenance-aware system for generating
instruction-tuning data for ESG and carbon-accounting tasks. The system
combines corpus processing, layered expert selection, multi-expert
collaboration, MetaExpert-guided refinement, and explicit source mapping
into a single reproducible data pipeline.
Our updated evaluation shows that CoDE substantially outperforms Direct
Prompting, Self-Instruct, and WizardLM-Evol under the finalized
protocol. On a 2,000-pair evaluation subset, CoDE achieves 4.0x higher
ROUGE-L, 4.6x higher BLEU-4, 10.5x higher Distinct-3, and 14.6x higher
domain relevance than the strongest baselines, while maintaining
FactCheck of 0.9478 and near-perfect structural completeness.
The VeritasCarbon-ESG-35K dataset demonstrates that domain-specific
instruction construction benefits from workflow control and provenance
preservation, not only from larger models or more elaborate prompts. The
finalized ablations further show that K = 3 is the best quality-cost
setting, collaboration improves quality over isolated experts, and
feedback remains necessary for viable CoDE generation. Future work
includes downstream fine-tuning experiments, stronger learned quality
evaluators, cross-document lineage tracking, and multilingual ESG
extensions.\(K\)
ACKNOWLEDGMENTS
The authors gratefully acknowledge the computing resources provided by
HKUST(GZ) HPC. We also thank the project collaborators who helped
complete the full VLDB experimental pipeline.
During the preparation of this work, the authors used generative AI
tools (including Kimi, ChatGPT-4o, and Claude 3.5) to assist with
language refinement, formatting, and the generation of supplementary
materials such as reproducibility checklists and prompt templates. All
factual claims, experimental results, and core technical contributions
were authored and verified by the human authors. The authors take full
responsibility for the content of this work.
\section{References}
\protect\phantomsection\label{ref_1}{}[1] Bender, E.M. and Friedman,
B. Data Statements for Natural Language Processing: Toward Mitigating
System Bias and Enabling Better Science. TACL, 2018.
\protect\phantomsection\label{ref_2}{}[2] Bingler, J.A. et al. Cheap
Talk and Cherry-Picking: What ClimateBert Has to Say on Corporate
Climate Risk Disclosures. Finance Research Letters, 2022.
\protect\phantomsection\label{ref_3}{}[3] Brown, T. et al. Language
Models are Few-Shot Learners. NeurIPS 2020.
\protect\phantomsection\label{ref_4}{}[4] Buneman, P. et al. Why and
Where: A Characterization of Data Provenance. ICDT 2001.
\protect\phantomsection\label{ref_5}{}[5] Chen, L. et al. AlpaGasus:
Training a Better Alpaca with Fewer Data. ICLR 2024.
\protect\phantomsection\label{ref_6}{}[6] Chen, W. et al.
AgentVerse: Facilitating Multi-Agent Collaboration and Exploring
Emergent Behaviors. ICLR 2024.
\protect\phantomsection\label{ref_7}{}[7] Cheney, J. et al.
Provenance in Databases: Why, How, and Where. Foundations and Trends in
Databases, 2009.
\protect\phantomsection\label{ref_8}{}[8] Chiang, W.-L. et al.
Vicuna: An Open-Source Chatbot Impressing GPT-4 with 90\% ChatGPT
Quality. 2023.
\protect\phantomsection\label{ref_9}{}[9] Chu, X. et al. Data
Cleaning: Overview and Emerging Challenges. SIGMOD 2016.
\protect\phantomsection\label{ref_10}{}[10] Cui, Y. et al. Tracing
the Lineage of View Data in a Warehousing Environment. ACM TODS, 2000.
\protect\phantomsection\label{ref_11}{}[11] Davidson, S.B. and
Freire, J. Provenance and Scientific Workflows: Challenges and
Opportunities. SIGMOD Record, 2008.
\protect\phantomsection\label{ref_12}{}[12] Doan, A. and Halevy, A.
Semantic Integration Research in the Database Community. AI Magazine,
2005.
\protect\phantomsection\label{ref_13}{}[13] Du, Y. et al. Improving
Factuality and Reasoning in Language Models through Multiagent Debate.
ICML 2024.
\protect\phantomsection\label{ref_14}{}[14] Friederich, D. et al.
Automated Identification of Climate Risk Disclosures in Annual Corporate
Reports. arXiv:2108.01415, 2021.
\protect\phantomsection\label{ref_15}{}[15] Gebru, T. et al.
Datasheets for Datasets. Communications of the ACM 64(12), 86-92, 2021.
\protect\phantomsection\label{ref_16}{}[16] Grafberger, S. et al.
Data Distribution Debugging in Machine Learning Pipelines. The VLDB
Journal 31(5), 1103-1126, 2022.
\protect\phantomsection\label{ref_17}{}[17] Green, T.J. et al.
Provenance Semirings. PODS 2007.
\protect\phantomsection\label{ref_18}{}[18] Herschel, M. et al. A
Survey on Provenance: What For? What Form? What From? The VLDB Journal
26(6), 881-906, 2017.
\protect\phantomsection\label{ref_19}{}[19] Hong, S. et al. MetaGPT:
Meta Programming for A Multi-Agent Collaborative Framework. ICLR 2024.
\protect\phantomsection\label{ref_20}{}[20] Honovich, O. et al.
Unnatural Instructions: Tuning Language Models with (Almost) No Human
Labor. ACL 2023.
\protect\phantomsection\label{ref_21}{}[21] Huang, A.H. et al.
FinBERT: A Large Language Model for Extracting Information from
Financial Text. Contemporary Accounting Research, 2023.
\protect\phantomsection\label{ref_22}{}[22] Ilyas, I.F. and Chu, X.
Data Cleaning. Morgan \& Claypool Publishers, 2019.
\protect\phantomsection\label{ref_23}{}[23] Kolbel, J.F. et al. Ask
BERT: How Regulatory Disclosure of Transition and Physical Climate Risks
Affects the CDS Term Structure. Journal of Financial Econometrics, 2022.
\protect\phantomsection\label{ref_24}{}[24] Lewis, P. et al.
Retrieval-Augmented Generation for Knowledge-Intensive NLP Tasks.
NeurIPS 2020.
\protect\phantomsection\label{ref_25}{}[25] Li, G. et al. CAMEL:
Communicative Agents for Mind Exploration of Large Language Model
Society. NeurIPS 2023.
\protect\phantomsection\label{ref_26}{}[26] Li, Y. et al.
Self-Alignment with Instruction Backtranslation. ICLR 2024.
[27] Lin, C.-Y. ROUGE: A Package for Automatic Evaluation of
Summaries. ACL Workshop, 2004.
\protect\phantomsection\label{ref_28}{}[28] Liu, W. et al. What
Makes Good Data for Alignment? A Comprehensive Study of Automatic Data
Selection in Instruction Tuning. ICLR 2024.
\protect\phantomsection\label{ref_29}{}[29] Luccioni, A.S. et al.
Estimating the Carbon Footprint of BLOOM, a 176B Parameter Language
Model. JMLR 2023.
\protect\phantomsection\label{ref_30}{}[30] Mitchell, M. et al.
Model Cards for Model Reporting. FAT* 2019.
\protect\phantomsection\label{ref_31}{}[31] Mukherjee, S. et al.
Orca: Progressive Learning from Complex Explanation Traces of GPT-4.
arXiv:2306.02707, 2023.
\protect\phantomsection\label{ref_32}{}[32] Ouyang, L. et al.
Training Language Models to Follow Instructions with Human Feedback.
NeurIPS 2022.
[33] Papineni, K. et al. BLEU: a Method for Automatic Evaluation of
Machine Translation. ACL, 2002.
\protect\phantomsection\label{ref_34}{}[34] Park, J.S. et al.
Generative Agents: Interactive Simulacra of Human Behavior. UIST 2023.
\protect\phantomsection\label{ref_35}{}[35] Polyzotis, N. et al.
Data Management Challenges in Production Machine Learning. SIGMOD 2017.
\protect\phantomsection\label{ref_36}{}[36] Rafailov, R. et al.
Direct Preference Optimization: Your Language Model is Secretly a Reward
Model. NeurIPS 2023.
\protect\phantomsection\label{ref_37}{}[37] Rahm, E. and Do, H.H.
Data Cleaning: Problems and Current Approaches. IEEE Data Eng. Bulletin,
2000.
\protect\phantomsection\label{ref_38}{}[38] Rekatsinas, T. et al.
HoloClean: Holistic Data Repairs with Probabilistic Inference. PVLDB
11(9), 2018.
\protect\phantomsection\label{ref_39}{}[39] Sanh, V. et al.
Multitask Prompted Training Enables Zero-Shot Task Generalization. ICLR
2022.
\protect\phantomsection\label{ref_40}{}[40] Schelter, S. et al.
Automatically Tracking Metadata and Provenance of Machine Learning
Experiments. NIPS 2017 Workshop on ML Systems.
\protect\phantomsection\label{ref_41}{}[41] Taori, R. et al.
Stanford Alpaca: An Instruction-Following LLaMA Model. Stanford
University, 2023.
\protect\phantomsection\label{ref_42}{}[42] Touvron, H. et al.
LLaMA: Open and Efficient Foundation Language Models. arXiv:2302.13971,
2023.
\protect\phantomsection\label{ref_43}{}[43] Touvron, H. et al. Llama
2: Open Foundation and Fine-Tuned Chat Models. arXiv:2307.09288, 2023.
\protect\phantomsection\label{ref_44}{}[44] Tunstall, L. et al.
Zephyr: Direct Distillation of LM Alignment. arXiv:2310.16944, 2023.
\protect\phantomsection\label{ref_45}{}[45] Varini, F.S. et al.
ClimaText: A Dataset for Climate Change Topic Detection. NeurIPS 2020
Workshop on Tackling Climate Change with ML.
\protect\phantomsection\label{ref_46}{}[46] Wang, Y. et al.
Self-Instruct: Aligning Language Models with Self-Generated
Instructions. ACL 2023.
\protect\phantomsection\label{ref_47}{}[47] Wang, Y. et al. How Far
Can Camels Go? Exploring the State of Instruction Tuning on Open
Resources. NeurIPS 2023 D\&B.
\protect\phantomsection\label{ref_48}{}[48] Wang, L. et al. A Survey
on Large Language Model based Autonomous Agents. Frontiers of Computer
Science, 2024.
\protect\phantomsection\label{ref_49}{}[49] Webersinke, N. et al.
ClimateBERT: A Pretrained Language Model for Climate-Related Text. arXiv
2021.
\protect\phantomsection\label{ref_50}{}[50] Wei, J. et al. Finetuned
Language Models are Zero-Shot Learners. ICLR 2022.
\protect\phantomsection\label{ref_51}{}[51] Wei, J. et al.
Chain-of-Thought Prompting Elicits Reasoning in Large Language Models.
NeurIPS 2022.
\protect\phantomsection\label{ref_52}{}[52] Whang, S.E. et al. Data
Collection and Quality Challenges in Deep Learning: A Data-Centric AI
Perspective. The VLDB Journal 32(4), 791-813, 2023.
\protect\phantomsection\label{ref_53}{}[53] Wu, S. et al.
BloombergGPT: A Large Language Model for Finance. arXiv 2023.
\protect\phantomsection\label{ref_54}{}[54] Wu, Q., Bansal, G.,
Zhang, J. et al. AutoGen: Enabling Next-Gen LLM Applications via
Multi-Agent Conversation. 2023.
\protect\phantomsection\label{ref_55}{}[55] Xie, Q. et al. PIXIU: A
Large Language Model, Instruction Data and Evaluation Benchmark for
Finance. arXiv:2306.05443, 2023.
\protect\phantomsection\label{ref_56}{}[56] Xu, C. et al. WizardLM:
Empowering Large Pre-Trained Language Models to Follow Complex
Instructions. ICLR 2024.
\protect\phantomsection\label{ref_57}{}[57] Xu, C. et al. Baize: An
Open-Source Chat Model with Parameter-Efficient Tuning on Self-Chat
Data. EMNLP 2023.
\protect\phantomsection\label{ref_58}{}[58] Yang, A. et al. Qwen2
Technical Report. arXiv 2024.
\protect\phantomsection\label{ref_59}{}[59] Yang, H. et al. FinGPT:
Open-Source Financial Large Language Models. IJCAI 2023 FinLLM Workshop.
\protect\phantomsection\label{ref_60}{}[60] Zha, D. et al.
Data-centric Artificial Intelligence: A Survey. ACM Computing Surveys,
2023.
\protect\phantomsection\label{ref_61}{}[61] Zhang, S. et al.
Instruction Tuning for Large Language Models: A Survey.
arXiv:2308.10792, 2023.
[62] Zhang, T. et al. BERTScore: Evaluating Text Generation with
BERT. ICLR 2020.
\protect\phantomsection\label{ref_63}{}[63] Zheng, L. et al. Judging
LLM-as-a-Judge with MT-Bench and Chatbot Arena. NeurIPS 2023.
\protect\phantomsection\label{ref_64}{}[64] Zhou, C. et al. LIMA:
Less Is More for Alignment. NeurIPS 2023.
\section{APPENDIX}
\textbf{Dataset as a Core Contribution}
In addition to the CoDE multi-agent framework, this paper makes a
substantial dataset contribution by releasing VeritasCarbon-ESG-35K, the
largest publicly available traceable ESG instruction dataset to date.
The dataset contains 35,009 high-quality instruction-response pairs
generated from 17,721 real-world ESG disclosure documents spanning four
corpus layers (domain textbooks, CSR reports, regulatory guidelines, and
industry analyses). Every pair is fully traceable to its source document
chunk and includes rich generation metadata (quality score, selected
experts, collaboration mode, knowledge injection count), enabling
fine-grained analysis of data provenance and quality.
Unlike existing ESG corpora that are limited to classification labels or
short extracts, VeritasCarbon-ESG-35K provides task-oriented
instructions with diverse types (Q\&A, summarization, extraction,
classification, analysis, temporal analysis, benchmark comparison,
greenwashing detection, standard alignment, knowledge graph
construction, and consistency verification). This makes it directly
usable for instruction tuning of large language models in the ESG and
sustainable-finance domains.
The dataset is released under a tiered strategy to balance accessibility
and copyright constraints:
\begin{itemize}
\item
  A representative 2,000-pair sample is included in the GitHub
  repository for immediate inspection and replication of the main
  evaluation (Table 2).
\item
  The full 35,009-pair dataset is hosted on Hugging Face Datasets
  (https://huggingface.co/datasets/Yihan-JIANG/VeritasCarbon-ESG-35K)
  under CC BY-SA 4.0.
\item
  Complete provenance metadata for all 17,721 source documents is
  provided in the repository, and the full dataset can be reproduced
  from source using the provided notebooks.
\end{itemize}
\textbf{A.1 Data and Code Availability}
\textbf{Code.} The complete VeritasCarbon framework, evaluation scripts,
and reproducibility instructions are available at
https://github.com/YihanJIANG-lab/VeritasCarbon-VLDB2027 under the MIT
License.
\textbf{Sample Data.} A representative 2,000-pair sample (8.7 MB) is
included in the GitHub repository. It is drawn with random.seed(42) from
the full 35,009-pair pool and matches the evaluation protocol used in
Table 2.
\textbf{Full Dataset.} The complete VeritasCarbon-ESG-35K dataset
(35,009 pairs, \textasciitilde153 MB) is hosted on Hugging Face
Datasets. It can be loaded in one line:
\begin{quote}
from datasets import load\_dataset
ds = load\_dataset("Yihan-JIANG/VeritasCarbon-ESG-35K", split="train")
\end{quote}
\textbf{Raw Corpus.} The source ESG documents (\textasciitilde2.7 GB,
17,721 files) contain copyrighted material and cannot be redistributed
in full. Complete provenance metadata is provided in
data/raw\_corpus/CORPUS\_MANIFEST.md.
\textbf{A.2 Reproducibility Checklist}
Table 1 provides a 30-minute verification checklist. All commands run on
CPU; no GPU is required.
Notes. (1) Tracks A-B require only the GitHub repo. (2) Track C requires
the raw corpus and an A100 GPU. (3) All random sampling uses
random.seed(42).
\textbf{A.3 Prompt Templates}
CoDE employs 11 specialized expert agents plus a MetaExpert. Each expert
receives a system prompt defining its persona, task, and output format.
All prompts support Default mode (generate from chunk) and Dynamic mode
(receive work instruction from MetaExpert). The dynamic mode prepends a
work instruction; the remainder is identical. Below we present the
default-mode prompts for representative experts, followed by a summary
table for all 11 experts, and finally the MetaExpert orchestration
prompt.
\textbf{A.3.1 Q\&A Expert}
\begin{quote}
You are a senior ESG (environment, social, governance) Q\&A generation
expert.
Task: Generate one high-quality Q\&A pair from the given text.
Requirements:
1. The question should be highly relevant to the text.
2. The answer must be findable in the text or reasonably inferred.
3. The question should cover an ESG dimension and lead to carbon
assessment.
4. The question should be specific and clear.
5. The answer should be accurate and complete.
\{knowledge\_context\}
Text:
\{chunk\_text\}
Output format:
Question: [question]
Answer: [answer]
Output only the Q\&A pair.
\end{quote}
\textbf{A.3.2 Analysis Expert}
\begin{quote}
You are a senior ESG analysis expert.
Task: Generate one analysis task from the text.
Requirements:
1. The instruction should ask to analyze ESG content (comparison, trend,
impact).
2. The answer should contain in-depth analysis and cross-domain
reasoning.
3. Be critical; assess credibility of carbon commitments.
\{knowledge\_context\}
Text:
\{chunk\_text\}
Output format:
Instruction: [instruction]
Analysis result: [analysis]
\end{quote}
\textbf{A.3.3 Greenwashing Detection Expert}
\begin{quote}
You are a senior ESG greenwashing detection expert.
Task: Generate a greenwashing detection task from the text.
Requirements:
1. Identify greenwashing patterns.
2. Quantify risk (0-1, 1 = highest risk).
3. Provide evidence chain.
Common patterns: vague commitments, selective disclosure,
exaggerated impact, time mismatch, scope confusion,
inconsistency, lack of verification.
\{knowledge\_context\}
Text:
\{chunk\_text\}
Output format:
Instruction: [detection instruction]
Answer: [matched patterns, risk level, evidence]
\end{quote}
\textbf{A.3.4 Consistency Verification Expert}
\begin{quote}
You are a senior ESG consistency verification expert.
Task: Generate a consistency verification task.
Requirements:
1. Check data consistency across the text.
2. Identify contradictions.
3. Detect missing-data patterns.
\{knowledge\_context\}
Text:
\{chunk\_text\}
Output format:
Instruction: [verification instruction]
Answer: [consistency assessment, contradictions]
\end{quote}
\textbf{A.3.5 Summary of All 11 Experts}
Table 2 summarizes each expert. All share the same prompt structure;
only task description and output labels differ.
\textbf{A.3.6 MetaExpert Orchestrator}
The MetaExpert designs work instructions for subordinate experts,
evaluates their outputs, and provides iterative feedback. Its prompt
dynamically incorporates the target expert type, core topics, chunk
preview, and collaboration context.
\begin{quote}
\# Role
You are a top carbon-strategy analyst.
Design one work instruction for a subordinate analyst.
\# Context
- Subordinate: \{expert\_type\} (\{expert\_desc\})
- Section: \{section\_name\}
- Core topics: \{topics\_str\}
- Preview: \{chunk\_text[:500]\}...
- Collaboration: experts=\{all\_experts\}
Current=\{expert\_type\}; Others=\{other\_experts\}
\# Task
Instruction must:
1. Carbon-centric: assess carbon performance or commitment credibility.
2. Cross-domain: use chunk info as evidence for carbon strategy.
3. Depth: compare, evaluate, predict, or critique-\/-not just extract.
\# Output
Work instruction:
\end{quote}
\textbf{Feedback round.} After receiving the expert\textquotesingle s
output, the MetaExpert scores correctness (0.4), completeness (0.3),
relevance (0.2), and structure (0.1). If any dimension falls below
threshold, a refined work instruction is generated and the expert
re-executes, up to R=2 rounds.
\textbf{A.4 Hyperparameters}
Table 3 lists all fixed hyperparameters. All random sampling uses
random.seed(42).
\textbf{A.5 Expert Selection Logic}
\begin{enumerate}
\def\labelenumi{\arabic{enumi}.}
\item
  Feature extraction. The chunk is analyzed for: keyword presence
  (carbon, emission, energy, employee, training, governance); numerical
  data density; temporal markers (years, quarters, since, by 20xx);
  standard references (GRI, TCFD, SASB, ISO 14001, CDP); comparative
  language; entity types (company names, plant locations, regulatory
  bodies).
\item
  Expert scoring. Each expert defines \(s_{i}(chunk)\) in [0,1].
  Examples: TemporalAnalysisExpert gets high score if temporal markers
  \textgreater= 2 and numerical trends present;
  GreenwashingDetectionExpert gets high score if vague-commitment
  keywords and selective-disclosure patterns detected;
  KnowledgeGraphExpert gets high score if entity count \textgreater= 5
  and relation verbs present.
\item
  Thresholding. Experts with \(s_{i} \geq 0.5\) are selected. If no
  expert meets the threshold, the fallback QAExpert is assigned.
\item
  Collaboration mode. If $\{s_i \geq 0.5\} \geq 2$,
  parallel collaboration is used; otherwise sequential.
\end{enumerate}
\textbf{MetaExpert override.} After expert selection, the MetaExpert may
dynamically adjust the assignment based on core topics. This override is
rare (\textless5\%) and logged in metadata.
\textbf{A.6 Error Analysis}
To illustrate the quality spectrum, we present representative examples
from the 2,000-pair sample. The quality score is a composite metric
(correctness 0.4, completeness 0.3, relevance 0.2, structure 0.1)
assigned by the MetaExpert.
\textbf{A.6.1 High-Quality Example (Score 0.915)}
\textbf{Experts:} Summary Expert + Consistency Verification Expert
\textbf{Instruction:} Generate a summary about the 2024-2025 A-share ESG
fund issuance trends, listed companies\textquotesingle{} information
disclosure quality, carbon-risk considerations in investment decisions,
corporate carbon-emission management progress, ESG evaluation-system
effectiveness, and the potential of China\textquotesingle s ESG market
to promote green transformation.
\textbf{Answer:} Since 2024, A-share ESG public fund issuance has
reached approximately 216 billion yuan. Although the issuance pace has
slowed, 18 new funds were still issued from Q4 to Q3 of the following
year. The industry\textquotesingle s overall scale has rebounded against
the trend, showing a structural shift from quantity expansion to quality
improvement. Under the dual-wheel drive of policy and market,
China\textquotesingle s ESG market is gradually maturing, with
carbon-risk considerations becoming an important factor in investment
decisions.
\textbf{Why it scores high:} The instruction demands multi-faceted
synthesis, and the answer integrates quantitative evidence (216B yuan,
18 funds) with qualitative assessment. The Consistency Verification
Expert cross-checks that fund-scale figures are internally coherent.
\textbf{A.6.2 Low-Quality Example (Score 0.000)}
\textbf{Experts:} Summary Expert
\textbf{Instruction:} Generate a summary for the following text.
\textbf{Answer:} This article aims to analyze how ESG information
disclosure affects employees\textquotesingle{} green innovation
efficiency and its role in promoting companies\textquotesingle{} real
carbon-reduction commitments. Special attention is paid to identifying
and criticizing potential greenwashing behavior, i.e., phenomena where
companies superficially promote environmental measures without taking
substantive action. In addition, the article...
\textbf{Why it scores low:} The output is a generic abstract of the
source text rather than a task-oriented instruction-response pair. The
instruction is merely generate a summary, and the answer restates the
paper\textquotesingle s purpose without concrete data or analysis. This
pattern occurs when the MetaExpert\textquotesingle s work instruction is
too vague and no second expert is assigned to refine depth.
\textbf{A.6.3 Key Failure Patterns}
Table 4 summarizes the most common failure modes observed in the
bottom-decile records.
\textbf{A.7 Dataset Schema}
Each record in the dataset is a JSON object with the following schema.
All fields are guaranteed present; optional sub-fields within metadata
default to null or empty lists when unavailable.
\textbf{Metadata Object}
\textbf{Quality score composition.} The MetaExpert evaluates four
dimensions, each scored 0-1: Correctness (weight 0.4), Completeness
(0.3), Relevance (0.2), and Structure (0.1). The composite score is the
weighted sum. Records with score \textless{} 0.6 trigger a feedback
round; records that still fail after R=2 rounds are discarded.
\textbf{A.8 Latency and Cost Breakdown}
Table 5 reports end-to-end latency measured on a single NVIDIA A100 (80
GB) with Qwen2-72B-Instruct 4-bit quantized via Unsloth. Batch size is
1; times are wall-clock averages over 100 records.
Notes. (1) Document parsing and chunking are I/O-bound and run on CPU.
(2) Generation and feedback are GPU-bound and dominate the pipeline. (3)
Parallel collaboration reduces per-record latency by invoking multiple
experts in a single forward pass when possible. (4) The 46-hour total
assumes sequential processing; with 4-way parallel chunk processing
across 4 A100 GPUs, the wall-clock time drops to \textasciitilde12
hours.
\end{document}
